\section{Technical Approach}
\label{sec:approach}
\begin{table}[t]\centering\footnotesize
\caption{Draft label-based era classes for our three target houses. Printed codes are the OCR
targets for weeks 6--7 (SS = spring/summer, AW = autumn/winter).}
\label{tab:eras}
\setlength{\tabcolsep}{3pt}
\begin{tabular}{@{}>{\raggedright\arraybackslash}p{1.25cm}>{\raggedright\arraybackslash}p{3.15cm}>{\raggedright\arraybackslash}p{3.15cm}@{}}\toprule
House & Era classes (label cue) & Printed date code \\\midrule
CDG & pre-1988 (no AD year); 1988--2000 (AD year, hyphenated style code); 2001--07; 2008--17;
2018-- (season-letter code cycles) &
``AD''~+ production year on the care label (SS goods usually carry the previous year); style-code season
letter A\,=\,SS2008\dots T\,=\,AW2017, restarting at A\,=\,SS2018 \\\addlinespace
Margiela & 1988--96 (blank label, four white tacks); 1997--2008 (numbered 0--23 label; women's Line~1 stays blank until c.\,2008);
2009--14 (design team, ``Maison Martin Margiela''); 2015-- (``Maison Margiela'') &
season label ``Printemps-\'Et\'e / Automne-Hiver'' + year (to c.\,1993); care-tag COMM.\ number
encoding year and season (1999--2010); ``YYYY 01/02'' from AW2010 \\\addlinespace
Undercover & 1990--2002 (Tokyo era); 2003--14 (Paris era, men's line sold as UNDERCOVERISM);
2015--20 (unified label); 2021-- &
product code on care/price tags in seven successive formats, e.g.\ ``UC''~+ season letter
P\,=\,AW2015\dots Z\,=\,AW2020; from SS2021 ``UC''~+ 1/2 (SS/AW) + year letter (A\,=\,2021) \\
\bottomrule\end{tabular}\end{table}

\begin{table*}[t]\centering\footnotesize
\caption{Proxy datasets used to validate the pipeline. Splits are grouped by product (merged when products share a photo), garment or runway show, and stratified by class.}\label{tab:datasets}
\begin{tabular}{lrlcllr}\toprule
Dataset & Images & Split groups & Brands & Eras & Train/val/test & License \\\midrule
Fashion-Gen & 18,967 & 3,916 (4,562 products) & 9 & 5 (years 2014--18) & 13,550/2,711/2,706 & research only \\
Second-hand labels & 2,321 & 2,321 garments & 11 & --- (---) & 1,657/332/332 & CC BY 4.0 \\
Europeana runway & 7,586 & 509 shows & 10 & 3 (decades 1990s--2010s) & 5,310/1,129/1,147 & In Copyright \\
\bottomrule\end{tabular}\end{table*}


\paragraph{Problem formulation.}
Given a photo $x$ of a second-hand garment (ideally of its label), we predict its brand
$b \in \mathcal{B}$ (the design house; sub-lines such as \emph{Homme Plus} or \emph{MM6} are folded
into the house), its era $e \in \mathcal{E}$, and its market value $v$. Brand and era are classification
heads on an image encoder $f(x)\in\mathbb{R}^d$, and value is estimated by retrieval from the prices
of the nearest labelled items in an embedding space (so far that of the frozen encoder).
On the real data, eras are defined per brand by label changes and printed date codes
(Sec.~\ref{sec:eras}); on the
proxy datasets they are whatever temporal label each dataset provides.

\paragraph{Datasets.}
Before collecting our own photos, we validate the pipeline on three public proxy datasets that
each cover part of the problem (Table~\ref{tab:datasets}):
(i)~\textbf{Fashion-Gen}~\cite{rostamzadeh2018fashiongen}: SSENSE catalog photos with brand,
season and retail price; we keep the clothing of nine archive houses (Comme des Gar\c{c}ons,
Maison Margiela, Undercover, Yohji Yamamoto, Junya Watanabe, Issey Miyake, Helmut Lang,
Raf Simons, Prada) and use the season year (2014--2018) as era and the list price (MSRP) as
value;
(ii)~\textbf{Second-hand labels}~\cite{secondhand2024dataset}: brand-label close-ups of donated
garments photographed at sorting stations, the closest public match to our in-store label photos
(mostly Nordic high-street brands, no era); and
(iii)~\textbf{Europeana runway}~\cite{europeana_api}: catwalk photos of ten archive designers with
show dates, giving real decade labels: 2000s and 2010s for all ten, plus late-1990s shows of CDG,
Yohji Yamamoto, Gaultier and Westwood (none for Margiela or Undercover).
Splits are grouped: all photos of one product, garment or show, and of products that share an
identical or near-identical photo (32$\times$32 grey thumbnails correlated at $\geq$0.99 and within
3 grey levels on average), fall in one split; they are stratified by class (5:1:1, i.e.\
71/14/14\%, overall). Classes with fewer than 50 images or 3 groups are dropped, and a class with
fewer than seven groups gets one group in validation, one in test and the rest in training, so every
kept class appears in each split. We share code, never images.

\paragraph{Frozen-feature baselines.}
We compare three pretrained encoders: CLIP ViT-B/16~\cite{radford2021clip}
(post-layernorm \texttt{[CLS]} token, 768-d), DINOv2 ViT-S/14~\cite{oquab2024dinov2}
(\texttt{[CLS]} concatenated with the mean patch token, 768-d, the paper's linear-evaluation
feature), and an ImageNet~\cite{deng2009imagenet} ResNet-50~\cite{he2016resnet} (2048-d).
On cached features we fit (a)~a linear probe: class-balanced multinomial logistic regression on
standardized features, with $C\in\{0.01,0.1,1,10\}$ chosen on validation macro-F1, following
the transfer protocol of~\cite{kornblith2019transfer}; and (b)~a $k$-NN classifier: cosine
similarity with votes weighted by $\exp(s/0.07)$, $k\in\{1,5,10,20,50\}$ chosen on validation.
$k$-NN doubles as a visual-search baseline in the spirit of retrieval-based fashion
recognition~\cite{liu2016deepfashion,kiapour2015where}.

\paragraph{Fine-tuning.}
The best frozen backbone is fine-tuned with one linear head per task and the sum of
class-balanced cross-entropies (label smoothing 0.1). The first epoch trains only the randomly
initialized heads on the frozen encoder while the learning rate warms up; later epochs unfreeze
everything with a 50$\times$ smaller encoder learning rate (AdamW~\cite{loshchilov2019adamw},
$10^{-3}$ heads / $2\cdot10^{-5}$ encoder, weight decay 0.05, cosine decay, 8 epochs, batch 64).
This only loosely follows linear-probe-then-fine-tune~\cite{kumar2022finetuning}, which distorts
pretrained features less and matters because our real test photos are out of distribution: after
one epoch the heads are far from a fitted probe. Starting the heads from the fitted linear probe is
implemented and is our next ablation. The epoch with the best mean validation macro-F1 is kept.

\paragraph{Augmentation for in-store photos.}
Training images are perturbed to look like hurried phone photos of a label: random resized
crops, rotation ($\pm25^\circ$), perspective warps, Gaussian blur, a synthetic specular
highlight (a soft bright ellipse blended over the image, mimicking glare from a flash or store
lighting), colour jitter, and JPEG re-compression. Horizontal flips are disabled for label
images, since mirrored text never occurs in practice. A \emph{light} variant (random resized
crop only, plus horizontal flips where allowed) serves as the ablation.

\paragraph{Value estimation.}
Value is predicted by retrieval, as a reseller would: the similarity-weighted geometric mean price
of the $k$ most similar training items ($k$ chosen on validation). We compare it with a ridge
regression on log price and with median-price-per-brand lookups using the true or the
predicted brand. On Fashion-Gen the target is retail MSRP, a proxy for resale value; resale
prices come from our own listing collection (Sec.~\ref{sec:changes}).

\paragraph{OCR.}
Labels print much of what we need as text: the brand, the sub-line and, for some houses, a
production-year mark or season code (Table~\ref{tab:eras}). Our first version runs
EasyOCR~\cite{easyocr} (CRAFT~\cite{baek2019craft} detection, CRNN~\cite{shi2017crnn} recognition),
letting it also try each text box rotated, and matches brand names within single text lines; a
confident match overrides the image model. In weeks 6--7 we add brand-specific parsers for printed
date codes, so that a confidently read year overrides the era head, and compare with
TrOCR~\cite{li2023trocr}.

\section{Evaluation Plan}
\label{sec:evalplan}
\textbf{Classification.} For brand and era we report accuracy, macro-F1 (classes are
imbalanced), balanced accuracy, top-3 accuracy, and row-normalized confusion matrices
(which eras are confused). Because the photos of one product are correlated, 95\% confidence
intervals come from a bootstrap that resamples whole products or shows. For the joint task we
report the fraction of images with both brand and era correct, and era accuracy among images
whose brand is correct. Hyperparameters are tuned only on validation; the test split is used
once per model.
\textbf{Value.} Median absolute percentage error (MdAPE), mean absolute error, the fraction of
items within $\pm25\%$ of the true price, and Spearman rank correlation.
\textbf{Robustness.} The fine-tuned models and the frozen linear probe of the same backbone (so far
CLIP) are scored on the test split under six corruptions (blur,
glare, rotation, perspective, JPEG compression, low light with sensor noise) at three
severities~\cite{hendrycks2019corruptions}, with identical seeded corruptions for all models.
\textbf{Ablations.} Backbone (CLIP vs.\ DINOv2 vs.\ ResNet-50), frozen vs.\ fine-tuned,
light vs.\ phone-photo augmentation, and image-only vs.\ image+OCR (done for brand on label photos;
for era once the date-code parsers exist).
\textbf{Real-world test.} The final test set is our own in-store phone photos of labels
(target: 20--30 garments per brand), never used for training or tuning; the gap between it and
the proxy test sets measures domain shift.
\textbf{Shortcut checks.} Era is also evaluated within each brand, since some brands appear
only in some years, and runway eras are taken from a single photo provider to limit how much photographic medium can
stand in for the decade; probes on low-level photo statistics measure what remains~\cite{geirhos2020shortcut}.
